\section{Opérateurs aréal et modulaire dans les secteurs \texorpdfstring{\(90\)}{90} et \texorpdfstring{\(120\)}{120}}

Dans chacune des deux familles de \cref{rm:eq:modular-boundary-counts}, il existe dix-huit canaux qutrit. Pour \(m\in\{90,120\}\), définissons
\begin{equation}
\label{rm:eq:modular-number-operators}
N_m=\sum_{e\in\partial A_m}N_e,
\qquad \Spec(N_e)=\{0,1,2\}.
\end{equation}
La sortie de Barbero--Immirzi est utilisée ici comme résultat antérieur, sans être redérivée. Avec la normalisation aréale construite dans \cref{rm:eq:modular-theta-clock-a0}, il vient
\begin{equation}
\label{rm:eq:modular-area-operator}
\boxed{
\frac{\widehat{\mathcal A}_{\rm phys}}{\ell_P^2}
=\gammaH a_0(N_{90}+N_{120}),}
\qquad
\gammaH a_0=0.0016755940749224991.
\end{equation}

\begin{theorem}[Spectre fini du quantum d’aire]
\label{rm:thm:modular-area-spectrum}
Sur les trente-six liens de la coupe,
\begin{equation}
\label{rm:eq:modular-total-number-spectrum}
\Spec(N_{90}+N_{120})=\{0,1,\ldots,72\}.
\end{equation}
La multiplicité de la valeur propre \(n\) est le coefficient de \(z^n\) dans \((1+z+z^2)^{36}\). Par conséquent,
\begin{equation}
\label{rm:eq:modular-area-spectrum}
\boxed{
\Spec\!\left(\frac{\widehat{\mathcal A}_{\rm phys}}{\ell_P^2}\right)
=\{n\gammaH a_0:0\leq n\leq72\}.}
\end{equation}
Les coupes compatibles prolongent la famille \(\mathcal A_n^{\rm prot}=n\gammaH a_0\ell_P^2\).
\end{theorem}

\begin{proof}
Chaque opérateur \(N_e\) a pour fonction génératrice spectrale \(1+z+z^2\). Les trente-six facteurs commutent et agissent sur des facteurs tensoriels distincts ; la fonction génératrice du total est donc \((1+z+z^2)^{36}\), dont les degrés parcourent tous les entiers de zéro à soixante-douze. La multiplication par le scalaire positif \(\gammaH a_0\ell_P^2\) donne \cref{rm:eq:modular-area-spectrum}.
\end{proof}

Il faut distinguer deux procédés de quantification. La loi de coupe \cref{rm:thm:modular-triadic-cut} compte quatre-vingt-une capacités unitaires et détermine \(81\log3\) ; l’opérateur \cref{rm:eq:modular-area-operator} quantifie les occupations de trente-six liens de réponse, jusqu’à soixante-douze trits de nombre. Ils partagent le même feuillet de bord, mais pas la même observable.

La torsion modulaire produit
\begin{equation}
\label{rm:eq:modular-delta-value}
\Dmod
=0.0001111970500599773249414566131933856\ldots,
\end{equation}
et le Hamiltonien de bord
\begin{equation}
\label{rm:eq:modular-hamiltonian}
\boxed{
\HHM^{(A)}=-\log\rho_A
=cI+\Dmod(90N_{90}+120N_{120}).}
\end{equation}
Le qualificatif \emph{holographique} ne désigne pas une analogie : l’incidence \(6\to8\) est transportée vers les ensembles de bord \(\partial A_{90}\) et \(\partial A_{120}\), et l’état réduit réside dans l’algèbre engendrée par leurs opérateurs nombre.

\begin{theorem}[Reconstruction conjointe de la géométrie et de la mémoire]
\label{rm:thm:modular-area-memory}
Soient
\begin{equation}
\label{rm:eq:modular-ND}
N=N_{90}+N_{120},
\qquad
D=90N_{90}+120N_{120}.
\end{equation}
Alors
\begin{equation}
\label{rm:eq:modular-inverse-channels}
\boxed{
N_{90}=4N-\frac{D}{30},
\qquad
N_{120}=\frac{D}{30}-3N.}
\end{equation}
En particulier, \([\widehat{\mathcal A}_{\rm phys},\HHM^{(A)}]=0\), et le couple \((\widehat{\mathcal A}_{\rm phys},\HHM^{(A)}-cI)\) reconstruit les deux secteurs d’occupation.
\end{theorem}

\begin{proof}
Les opérateurs \(N_{90}\) et \(N_{120}\) agissent sur des facteurs distincts et commutent. L’application linéaire qui envoie \((N_{90},N_{120})\) sur \((N,D)\) a pour matrice
\[
\begin{pmatrix}1&1\\90&120\end{pmatrix},
\qquad \det=30.
\]
Son inverse est exactement \cref{rm:eq:modular-inverse-channels}. Puisque \cref{rm:eq:modular-area-operator,rm:eq:modular-hamiltonian} sont des combinaisons linéaires des mêmes opérateurs nombre, le commutateur s’annule.
\end{proof}

L’interprétation opératorielle est précise : l’observable d’aire conserve le nombre total d’excitations, tandis que le Hamiltonien modulaire conserve leur secteur d’origine. Aucune des deux observables ne retient à elle seule les deux coordonnées ; leur couple conjoint les détermine.

\section{Confluence de l’aire, de l’information et de l’énergie}

L’échelle angulaire \(\theta_{\rm clk}\) de \cref{rm:eq:modular-theta-clock-a0} et la température chronologique \(\Theta_{\rm clk}\) sont des objets distincts ; l’emploi de la majuscule fait partie du typage. La réalisation thermodynamique de \cref{rm:ch:metrology} démontre
\begin{equation}
\label{rm:eq:modular-planck-trit-bridge}
E_P=k_B\Theta_{\rm clk}\log3,
\end{equation}
tandis que la réalisation aréale construit l’incrément élémentaire
\begin{equation}
\label{rm:eq:modular-area-quantum}
\delta\mathcal A
=\gammaH a_0\ell_P^2.
\end{equation}
Elles ne sont pas égalisées : \cref{rm:eq:modular-planck-trit-bridge} fixe le quantum informationnel d’énergie et \cref{rm:eq:modular-area-quantum} le quantum géométrique d’aire. Leur quotient définit la tension de conversion typée
\begin{equation}
\label{rm:eq:modular-information-area-tension}
\boxed{
\mathcal T_{I/A}
=\frac{E_P}{\delta\mathcal A}
=\frac{k_B\Theta_{\rm clk}\log3}
{\gammaH a_0\ell_P^2}.}
\end{equation}
L’expression n’introduit pas de constante supplémentaire : elle détermine les normalisations qui doivent être conservées lors de la transformation de l’information en géométrie.


Dans la coupe minimale de quatre-vingt-un liens, la capacité informationnelle physique et l’énergie d’un trit par lien sont
\begin{equation}
\label{rm:eq:modular-cut-info-energy}
S_{\rm cut}^{\rm phys}=81k_B\log3,
\qquad
E_{\rm cut}^{(1)}=81E_P
=\Theta_{\rm clk}S_{\rm cut}^{\rm phys}.
\end{equation}
L’égalité finale est une équation d’état de cette réalisation, non une loi universelle d’équipartition. En particulier, l’état de Gibbs des trente-six liens de bord possède une entropie différente et est traité ci-dessous.

Cette séparation détermine le contenu structurel de la confluence : Barbero--Immirzi oriente le spectre d’aire ; \(\log3\) mesure la décision irréductible du TRIT ; \(\Theta_{\rm clk}\) convertit cette information en énergie ; et \(\ell_P^2\) relie le quantum d’aire à la famille gravitationnelle. La gravité est décrite comme confluence de réalisations fonctionnelles d’un antécédent commun, et non comme une égalité imposée entre leurs valeurs.
